where @@M0@@ is the @@M1@@-scheme
obtained by blowing up @@M2@@. We suppose in addition that @@M3@@ is quasi-compact and has an ample sheaf.

\sgaStableTarget{sga6.e08.prop.4-2}\textbf{Proposition 4.2.} \uline{The morphism} @@M4@@ \uline{is projective}, \uline{a complete intersection}, \uline{and of virtual relative dimension zero}. \uline{We also have the relation}:
@@M5@@

The first assertion follows immediately from (VII 1.8).

The second relation follows from the relation:
@@M6@@
proved in (VII 3.6).

\sgaStableTarget{sga6.e08.prop.4-3}\textbf{Proposition 4.3.} \uline{For every} @@M7@@, \uline{we have the relation}:
@@M8@@
\uline{where} @@M9@@ \uline{denotes the} @@M10@@-\uline{th Chern class}.

This relation is contained in VII 4.8, since it is precisely @@M11@@.

\sgaStableTarget{sga6.e08.prop.4-4}\textbf{Proposition 4.4.} \uline{Let} @@M12@@ \uline{be a quasi-compact scheme}. \uline{If} @@M13@@ \uline{is a locally free sheaf of finite type on} @@M14@@, \uline{of class} @@M15@@ \uline{in} @@M16@@, \uline{and of rank} @@M17@@, \uline{we have}:
@@M18@@

The two sides of the relation to be proved are multiplicative: this is clear for @@M19@@, and for @@M20@@; it also holds for @@M21@@, which is the class of @@M22@@. Consequently, by passing to the flag bundle of @@M23@@, one may suppose that @@M24@@ is a sum of classes of invertible sheaves, and reduce to proving the formula for an invertible sheaf. Let
